\chapter{Introducción}

\section{Contexto y motivación}

El aprendizaje automático (Machine Learning) está en su mejor momento. Hoy en día lo usamos para resolver problemas en la medicina, las finanzas o la ingeniería con resultados increíbles. Sin embargo, quienes trabajamos en esto nos enfrentamos a un dolor de cabeza constante: los datos del mundo real rara vez están completos.\parencite{little2019}.

Ya sea por el fallo de un sensor, un error humano al registrar la información, o la propia naturaleza del proceso, siempre nos topamos con ``huecos'' en nuestros datasets. Esto es un problema grave porque la gran mayoría de los algoritmos no saben qué hacer cuando les falta información. Antes de poder entrenar cualquier modelo, tenemos que tomar una decisión sobre cómo rellenar esos vacíos; un proceso que conocemos como imputación.\parencite{garcia2017}.

Y aquí está el truco: cómo rellenamos esos datos importa, y mucho. Si simplemente rellenamos los huecos con el promedio general (imputación por media), es rápido y fácil, pero podemos distorsionar la realidad de nuestros datos. Por otro lado, si usamos algoritmos más inteligentes que buscan patrones (como KNN, que se fija en datos ``vecinos'' similares), conservamos mejor la estructura de la información, pero nos cuesta más tiempo y recursos computacionales.\parencite{troyanskaya2001}.


\section{Problema de investigación}

A pesar de que la imputación de datos faltantes es una tarea fundamental
en el preprocesamiento de cualquier pipeline de ciencia de datos, existe
escasa evidencia empírica que evalúe de forma sistemática el impacto que
tiene la elección del método de imputación sobre el rendimiento final del
modelo de clasificación \parencite{salgado2016}. La mayoría de los estudios
existentes analizan los métodos de imputación de manera aislada, sin evaluar
el impacto integral de dicha elección sobre el rendimiento del clasificador
en el pipeline completo.

En este contexto, el presente trabajo plantea la siguiente pregunta de
investigación: ¿en qué medida el método de imputación utilizado —KNN frente
a la imputación por media— afecta el rendimiento de un modelo de
clasificación supervisada, medido mediante las métricas Accuracy y F1-score?

\section{Objetivos}

\subsection{Objetivo general}

Evaluar el impacto de métodos de imputación de datos faltantes —KNN frente
a métodos tradicionales como la imputación por media— sobre el rendimiento
de un modelo de clasificación supervisada, mediante un experimento controlado
y reproducible.

\subsection{Objetivos específicos}

\begin{itemize}
    \item Analizar las características estadísticas del \textit{Water
          Potability Dataset} mediante técnicas de análisis exploratorio
          de datos, identificando distribuciones, valores atípicos y
          valores faltantes preexistentes.

    \item Diseñar un experimento controlado que introduzca valores faltantes
          de forma aleatoria bajo el mecanismo MCAR (\textit{Missing
          Completely At Random}), con el fin de simular escenarios de
          datos incompletos en condiciones reproducibles.

    \item Implementar un pipeline automatizado en Python que aplique los
          métodos de imputación por media e imputación KNN de forma
          sistemática y reproducible.

    \item Comparar el efecto de cada método de imputación sobre las
          métricas Accuracy y F1-score de un clasificador de Regresión
          Logística entrenado sobre los datos imputados.

    \item Analizar las implicancias prácticas de los resultados obtenidos
          para el diseño de pipelines de ciencia de datos en escenarios
          con datos incompletos.
\end{itemize}

\section{Hipótesis}

Se plantea la hipótesis de que el método de imputación KNN produce un
mejor rendimiento en el modelo de clasificación en comparación con la
imputación por media. Esto se fundamenta en que KNN estima los valores
faltantes considerando la similitud entre instancias del dataset, lo que
permite preservar en mayor medida la estructura estadística y las
correlaciones originales entre variables, proporcionando así información
de mayor calidad al clasificador \parencite{troyanskaya2001}.

\section{Alcance y limitaciones}

El presente trabajo se limita a la comparación de dos métodos de imputación
—media y KNN— aplicados sobre un único dataset y evaluados mediante un
único modelo de clasificación base (Regresión Logística). No se aborda la
optimización de hiperparámetros ni la comparación entre múltiples
clasificadores, aspectos que quedan propuestos como líneas de trabajo futuro.

Asimismo, el experimento considera exclusivamente el mecanismo de ausencia
MCAR (\textit{Missing Completely At Random}), que asume que la probabilidad
de que un valor sea faltante es independiente de cualquier variable observada
o no observada \parencite{little2019}. Los mecanismos MAR (\textit{Missing
At Random}) y MNAR (\textit{Missing Not At Random}), que representan
escenarios de mayor complejidad, quedan fuera del alcance de este estudio.

